\section{The all-integer resonant evaluation}\label{sec:resonances}

\subsection{Reference shifts and a coefficient differential identity}
The coefficient polynomials have more structure than their recursive
definition initially suggests.

\begin{lemma}[Exact transport in the reference shift]\label{lem:shift}
For arbitrary reference shifts $c,d$,
\begin{equation}\label{eq:shift}
 A_j(u;c)=\sum_{l=0}^j A_l(u;d)
 \frac{(1+u+l)_{j-l}}{(j-l)!}(c-d)^{j-l}.
\end{equation}
In particular, $A_j$ has degree at most $j$ in $c$, and
\begin{equation}\label{eq:cderiv}
 \partial_c A_j(u;c)=(u+j)A_{j-1}(u;c),\qquad j\ge1.
\end{equation}
\end{lemma}
\begin{proof}
In the expansion with reference $d$, write
$n+d=(n+c)-(c-d)$ and expand each factor by the binomial theorem.
Uniqueness of asymptotic coefficients gives \eqref{eq:shift}.
Taking its first-order variation in $c-d$ proves \eqref{eq:cderiv}.
Since these are polynomial identities, establishing them for positive
references proves their algebraic form without further restrictions.
\end{proof}

\begin{lemma}[Coefficient differential law]\label{lem:parameterlaw}
For $j\ge1$,
\begin{equation}\label{eq:parameterlaw}
 (\partial_a+\partial_b-\partial_u)A_j
 =-(u+j)A_{j-1}+
 \sum_{k=1}^j\zeta(1-k,c)A_{j-k}.
\end{equation}
\end{lemma}
\begin{proof}
Treating $n$ temporarily as a positive real variable, logarithmic
differentiation of \eqref{eq:rn} gives the exact identity
\begin{equation}\label{eq:termparameter}
 (\partial_a+\partial_b-\partial_u)r_n(u)
 =\partial_n r_n(u)+\psi(n+1)r_n(u).
\end{equation}
The digamma asymptotic expansion at the same reference is
\begin{equation}\label{eq:digammaasymp}
 \psi(n+1)\sim\log(n+c)+
 \sum_{k\ge1}\frac{\zeta(1-k,c)}{(n+c)^k}.
\end{equation}
Indeed, its coefficients are $-B_k(c)/k=\zeta(1-k,c)$, using Bernoulli
reflection.  Substitute the finite expansions into \eqref{eq:termparameter}.
The logarithmic terms cancel because $-\partial_u(n+c)^{-1-u}$
is $\log(n+c)(n+c)^{-1-u}$.  At inverse-power order $j$,
$\partial_n r_n$ contributes $-(u+j)A_{j-1}$ and the nonlogarithmic
digamma terms contribute the sum in \eqref{eq:parameterlaw}.
Uniqueness of the asymptotic expansion completes the proof.
\end{proof}

\subsection{Residue cancellation at every nonpositive integer}
Set
\begin{equation}\label{eq:rho}
 \rho_N(a,b)=\frac{(-1)^N(1-a)_N(1-b)_N}{N!}.
\end{equation}
This is a polynomial, including when one of the factors vanishes.

\begin{proposition}[Two exact resonant coefficient identities]\label{prop:rho}
For every $N\ge0$,
\begin{align}
 A_N(-N;c)&=\rho_N(a,b),\label{eq:residueid}\\
 \partial_u A_N(-N;c)
 +\sum_{j=0}^{N-1}A_j(-N;c)\zeta(1-N+j,c)
 &=(\partial_a+\partial_b)\rho_N(a,b).
 \label{eq:bernoulliresonance}
\end{align}
The sum in the second identity is empty for $N=0$.
\end{proposition}
\begin{proof}
Take $K=N+1$ in Theorem~\ref{thm:subtraction}.
The only zeta factor with a pole at $u=-N$ has index $j=N$ and residue
$A_N(-N;c)$.  On the other hand,
$\Res_{u=-N}G(u)=\rho_N$, by the gamma residue formula and
$\Gamma(a)/\Gamma(a-N)=(a-N)_N=(-1)^N(1-a)_N$.
Since $R_{N+1}$ is holomorphic, its residue vanishes, proving
\eqref{eq:residueid}.  For generic parameters this is a direct residue
calculation; continuity gives the polynomial statement at degeneracies.

For $N\ge1$, put $j=N$, $u=-N$ in \eqref{eq:parameterlaw}.
The factor $u+j$ vanishes.  Use \eqref{eq:residueid} for the
$a,b$ derivatives and rearrange.  This proves
\eqref{eq:bernoulliresonance}.  For $N=0$, both sides are zero because
$A_0=\rho_0=1$.
\end{proof}

The second identity is an explicit finite Bernoulli-polynomial identity,
not merely a residue statement.  It is the cancellation needed to turn a
long finite-part formula into a short evaluation.

\begin{theorem}[Closed evaluation at every resonance]\label{thm:resonant}
For positive $a,b,c$ and every $N\ge0$,
\begin{align}
 &\sum_{n=0}^\infty\left\{
 \frac{\Gamma(n+a)\Gamma(n+b)}
      {\Gamma(n+1)\Gamma(n+a+b-N)}
 -\sum_{j=0}^{N}A_j(-N;c)(n+c)^{N-1-j}\right\}\notag\\
 &\hspace{15mm}=
 \boxed{\rho_N(a,b)
 \bigl(H_N-\gamma-\psi(a)-\psi(b)+\psi(c)\bigr).}
 \label{eq:closedresonance}
\end{align}
The series is absolutely convergent, including at positive integer
parameters for which $\rho_N=0$.
\end{theorem}
\begin{proof}
The case of nondegenerate parameters is enough initially.  Put $u=-N+t$.
The gamma Laurent expansion gives
\[
 G(-N+t)=\frac{\rho_N}{t}+
 \rho_N\{\psi(N+1)-\psi(a-N)-\psi(b-N)\}+O(t).
\]
Writing $A_N(-N+t;c)=\rho_N+t\partial_uA_N(-N;c)+O(t^2)$
and using \eqref{eq:stieltjes}, its matching zeta product has constant term
$\partial_uA_N(-N;c)-\rho_N\psi(c)$.
Consequently \eqref{eq:master} gives
\begin{align*}
 R_{N+1}(-N;c)
 ={}&\rho_N\{\psi(N+1)-\psi(a-N)-\psi(b-N)+\psi(c)\}\\
 &-\partial_uA_N(-N;c)
 -\sum_{j=0}^{N-1}A_j(-N;c)\zeta(1-N+j,c).
\end{align*}
The last line is $-(\partial_a+\partial_b)\rho_N$ by
\eqref{eq:bernoulliresonance}.  Also
\[
 \frac{\partial_a\rho_N}{\rho_N}=
 \psi(a)-\psi(a-N),\qquad
 \frac{\partial_b\rho_N}{\rho_N}=
 \psi(b)-\psi(b-N).
\]
Substitution cancels the shifted digammas, leaving
$\rho_N\{\psi(N+1)-\psi(a)-\psi(b)+\psi(c)\}$.
Finally $\psi(N+1)=H_N-\gamma$.
Both the normally convergent bracketed sum and this last expression are
continuous, indeed locally holomorphic, in the positive parameters.
The identity therefore extends across all zeros of $\rho_N$.
\end{proof}

\begin{remark}[What is classical, and what the proof adds]
The value theorem is consistent with the classical hypergeometric
constant-term analysis discussed in Section~\ref{sec:scope}.  The proof
above explains its simplification entirely within the fixed-shift
coefficient algebra, supplies an exact recurrence for every subtraction,
and remains valid at reciprocal-gamma degeneracies.  The next section
retains the entire Taylor expansion rather than only its constant term.
\end{remark}
